% =====================================================================
\section{A working pipeline}
\label{sec:pipeline}
% =====================================================================

The address of a declaration is computed in three stages:
canonicalization, serialization, and hashing.

% --------------------------------------------------------------------
\subsection{Canonicalization}
\label{sec:canon}
% --------------------------------------------------------------------

Canonicalization transforms a Lean~4 \texttt{Expr} (expression) into a
canonical form where expressions differing only in irrelevant ways
become syntactically identical.  We implement two levels, offering a
trade-off between stability and the strength of the equivalence
relation.

\paragraph{Level~0: structural canonicalization.}
Level~0 is a pure transformation (no monadic context required) that
normalizes two sources of superficial variation.  First, \emph{universe
parameters}: Lean uses named universe polymorphism, and these names are
arbitrary; Level~0 renames them to positional indices (\texttt{u\_0},
\texttt{u\_1}, \ldots) by order of first occurrence.  Second,
\emph{metadata annotations}: Lean attaches \texttt{.mdata} nodes to
expressions for internal bookkeeping (source positions, elaboration
hints); these carry no mathematical content and are stripped.

Lean~4 uses de~Bruijn indices~\cite{debruijn1972} for bound variables,
so $\alpha$-equivalence is handled by the representation itself.
Binder \emph{names} (e.g., \texttt{n} vs.\ \texttt{x}) are cosmetic
and excluded from hashing.  Level~0 is fast and deterministic.  It
depends only on the structure of the \texttt{Expr} type, which is
stable within a major Lean version but has changed across major versions
(e.g., Lean~3 to Lean~4).  Its limitation is that it does not normalize
definitional equalities: two formulations that differ by an unfolded
abbreviation or a reduced beta-redex receive different hashes.

\paragraph{Level~1: reducible normalization.}
Level~1 extends Level~0 with \emph{reducible-transparency
normalization}.  It requires a \texttt{MetaM} (meta-programming monad)
context because it invokes Lean's \texttt{whnf} (weak head normal form)
reduction.  The normalizer, \texttt{deepNormalize}, recursively applies
\texttt{whnf} at every subexpression under \texttt{reducible}
transparency:

\begin{lstlisting}[style=pseudo,caption={Deep normalization (simplified pseudocode)}]
deepNormalize(e, depth):
  if depth >= 256: return e       -- recursion guard
  e <- whnf e                     -- reducible transparency
  match e:
    .app f a      -> .app (deepNormalize f) (deepNormalize a)
    .lam _ t b _  -> normalize t; open binder; normalize body;
                     re-abstract
    .forallE _ t b _ -> same as lambda
    .letE _ _ v b -> normalize v; substitute; normalize result
    .mdata _ e    -> deepNormalize e
    .proj _ _ e   -> normalize e; whnf the projection
    _             -> return e
\end{lstlisting}

Under reducible transparency, notation aliases (e.g.,
\texttt{HAdd.hAdd} $\to$ \texttt{Add.add}), declarations marked
\texttt{@[reducible]}, beta-redexes, let-bindings, and reducible
structure projections are all unfolded.  Crucially,
\texttt{@[semireducible]} declarations (including typeclass instances)
are \emph{not} unfolded, preserving abstraction boundaries:
\texttt{Add Nat} stays as \texttt{Add Nat} rather than being replaced
by its instance implementation.

As an example, consider \texttt{theorem add\_comm\_v1 : $\forall$ (a b
: Nat), a + b = b + a} and a variant \texttt{add\_comm\_v2 :
$\forall$ (a b : Nat), b + a = a + b}.  At Level~0, these may receive
different hashes if the elaborated expressions differ structurally.  At
Level~1, after \texttt{whnf} normalization, both reduce to the same
canonical form and receive the same hash.  On
\texttt{Init.Data.Nat.Basic}, the largest equivalence class grows from
391 members (Level~0) to 467 (Level~1)---76~additional declarations
collapsed by reducible unfolding.  Section~\ref{sec:results} presents a
larger-scale evaluation.

\franz{How useful is Level~1 in practice? Probably handles ``boring'' statements, but two mathematicians independently formalizing the same interesting theorem---would they get the same hash? Almost never at Level~1.}

\paragraph{Level~2: full definitional equality (open problem).}
The ideal canonicalization would assign the same hash to all
definitionally equal types.  However, Lean's dependent type theory does
not admit unique normal forms under full definitional
equality~\cite{abel2018decidability}.  Achieving Level~2 would require
either restricting the fragment of the type theory considered, or
developing new normalization techniques.  We leave this as an open
problem (see Section~\ref{sec:equiv} for further discussion).

\franz{Even Level~2 may not suffice---people use slightly different formalizations.}

\franz{Propose Level~3: formally different but ``spiritually the same''---one trivially proves the other. Quantify with fuzzy match + trivial proof detection (e.g.\ \texttt{exact?}\ or \texttt{apply?}\ in one step).}

% --------------------------------------------------------------------
\subsection{Serialization, Hashing, and Hashing Modes}
\label{sec:serial}
% --------------------------------------------------------------------

\paragraph{Serialization.}
The canonical expression is serialized to a \texttt{ByteArray} using a
tagged encoding: each \texttt{Expr} constructor is assigned a unique tag
byte, followed by the serialization of its children.  Binder names are
omitted (de~Bruijn indices suffice).  \texttt{BinderInfo} (implicit,
explicit, instance, etc.)\ is included for lambda and forall nodes.
\texttt{.mdata} is transparent: the child is serialized directly.  For
\texttt{.const name levels}, the treatment of \texttt{name} depends on
the hashing mode (below).

\sam{Hash-based proof compression: replace subexpressions by hash on network layer. Sam's ``hexpr'' constructor elides subterm structure for storage/communication.}

\paragraph{Hashing.}
The serialized byte array is hashed with SHA-256, producing a
64-character hexadecimal string.  We use SHA-256 via OpenSSL's EVP API
through a minimal C FFI shim, chosen for ubiquity and ease of binding
rather than performance (a faster hash such as BLAKE3 would work equally
well---the choice of hash function is orthogonal to the addressing
design).
\simone{Address: (a) mutual recursion in Merkle DAG, (b) why SHA-256 not BLAKE3, (c) literal \texttt{PROOF\_IRRELEVANT} sentinel value, (d) is L0 truly Lean-version-independent given \texttt{Expr} changes?}

\paragraph{Hashing modes.}
\label{sec:modes}
When serializing a constant reference (\texttt{.const name levels}),
there is a choice of what to emit for \texttt{name}.  In
\emph{name-based} mode (default), the declaration's \texttt{Name} is
serialized as a UTF-8 string---fast and stable within a single library,
but if a dependency is renamed, the address changes.  In
\emph{content-based} (Merkle DAG) mode, the 32-byte content hash of the
referenced declaration is emitted instead.  Declarations are processed
in topological (dependency) order via Kahn's
algorithm~\cite{kahn1962topological}, so each hash incorporates the
hashes of its transitive dependencies.  The result is a Merkle
DAG~\cite{merkle1988} where a declaration's identity depends only on the
mathematical content of everything it (transitively) references.

Mutually recursive declarations form cycles in the dependency graph.
The topological sort detects these and falls back to name-based hashing
for cycle members---a conservative choice that preserves soundness at
the cost of weaker identity for mutual definitions.  Mutual recursion
is uncommon in theorem \emph{statements}, but Lean's elaboration of
inductive types generates mutually recursive auxiliary declarations
(\texttt{rec}, \texttt{casesOn}, etc.), so the fallback affects more
declarations than one might initially expect when indexing definitions
alongside theorems.

\begin{table}[h]
\centering
\begin{tabular}{@{}lll@{}}
  \toprule
  \textbf{Mode} & \textbf{Constant references become} & \textbf{Trade-off} \\
  \midrule
  Name-based & The declaration's \texttt{Name} string & Fast; changes if deps renamed \\
  Content-based & 32-byte content hash of the dep & True content identity; slower \\
  \bottomrule
\end{tabular}
\caption{Hashing modes.}
\label{tab:modes}
\end{table}


% --------------------------------------------------------------------
\subsection{Results}
\label{sec:results}
% --------------------------------------------------------------------

This section presents both concrete examples of content-addressing in
action and a quantitative evaluation on Mathlib.

\paragraph{Duplicate detection.}
Large libraries inevitably accumulate duplicates---theorems with
different names but identical mathematical content.  Content-addressing
detects these automatically: duplicates share the same hash.  The
technique of comparing type expressions for deduplication is already
used in the AI-for-theorem-proving community---for instance, synthetic
theorem generation pipelines deduplicate candidates by \texttt{Expr}
equality~\cite{syntheticthm}.  Content-addressing provides a
principled, persistent, and cross-project version of what is currently
done ad~hoc within individual pipelines.

\paragraph{Cross-project references.}
Consider two independent projects that both work with a mathematical
concept:

\begin{lstlisting}[caption={Two projects proving the same lemma independently}]
-- In project-a/Algebra/Basic.lean
theorem add_comm_nat : forall a b : Nat, a + b = b + a :=
  Nat.add_comm

-- In project-b/MyMath/Nat.lean
theorem nat_addition_commutative : forall a b : Nat, a + b = b + a :=
  fun a b => Nat.add_comm a b
\end{lstlisting}

Despite different names, files, and projects, both declarations receive
the same content address.  A search tool can discover that these are the
same mathematical fact, enabling cross-project deduplication and
discovery without any coordination between the projects.

\paragraph{Open problems and conditional results.}
Content-addressing supports a structured workflow for open problems,
though this aspect of the system is prospective---the mechanism is
implemented but has not yet been exercised at scale.  A researcher
publishes a conjecture as an \emph{open point}---a \texttt{Prop} with a
content address but no proof:

\begin{lstlisting}[caption={An open conjecture and a conditional result}]
@[open_point "Bertrand's postulate, elementary form"]
def BertrandPostulate : Prop :=
  forall n : Nat, n >= 1 -> exists p, Nat.Prime p /\ n < p /\ p < 2 * n

@[publish "Consequence of Bertrand's postulate"]
theorem prime_between (n : Nat) (h : BertrandPostulate) (hn : n >= 1) :
    exists p, Nat.Prime p /\ n < p /\ p < 2 * n :=
  h n hn
\end{lstlisting}

The registry (Section~\ref{sec:registry}) tracks which conditional
results depend on which open points.  When someone eventually proves the
conjecture, the status of all downstream results cascades from
``conditional'' to ``proved.''  Several practical questions remain open:
how to handle conjectures that turn out to be false (or independent),
how to distinguish \texttt{sorry}-based development from genuinely
axiomatic assumptions, and how to curate status metadata across
projects.  We discuss these further in Section~\ref{sec:future}.
\simone{Mark open-problem workflows (Sections 5.3--5.4) as speculative/future work. No real usage evidence. Address: who curates status? What if open point is false? \texttt{sorry}-based vs axiomatic?}

\paragraph{Registry generation.}
A project publishes its annotated declarations by adding
\texttt{\#ca\_registry "registry/"} at the end of a Lean file that
imports the annotated modules.  When \texttt{lake build} compiles this
file, it writes \texttt{declarations.json} (address, name, status,
type, deps for each declaration) and \texttt{meta.json} (project
summary with counts by status) to the specified directory.

\paragraph{Evaluation on Mathlib.}
We evaluated the pipeline on the environment loaded by
\texttt{Mathlib.\allowbreak Algebra.\allowbreak Group.\allowbreak Basic}, a core module of Mathlib that
transitively imports 1{,}209 modules covering foundational arithmetic,
data structures, and algebraic hierarchy definitions.  The pipeline was
run in name-based mode at Level~1 (reducible normalization), restricted
to theorems.  All 57{,}613 theorems in the transitive environment were
processed successfully---no declarations were skipped or failed.

The 57{,}613 theorems produce 57{,}036 unique type addresses.  Of
these, 544 addresses are shared by more than one theorem, covering
1{,}121 theorems (1.9\% of the total).

\begin{table}[h]
\centering
\begin{tabular}{@{}rrr@{}}
  \toprule
  \textbf{Class size} & \textbf{Classes} & \textbf{Theorems} \\
  \midrule
  1 (unique)  & 56{,}492 & 56{,}492 \\
  2           & 513      & 1{,}026 \\
  3           & 29       & 87 \\
  4           & 2        & 8 \\
  \midrule
  Total       & 57{,}036 & 57{,}613 \\
  \bottomrule
\end{tabular}
\caption{Equivalence class size distribution (L1, name-based, theorems only).}
\label{tab:eval}
\end{table}

\paragraph{Examples of detected duplicates.}
The equivalence classes reveal genuine duplicates---theorems that exist
under multiple names due to historical renaming, namespace aliasing, or
independent re-statement:

\begin{lstlisting}[style=output,caption={Selected equivalence classes (L1)}]
Size 4:
  by_contra = by_contradiction
           = Classical.byContradiction = Classical.by_contradiction
  Type: (Not p -> False) -> p

Size 3:
  congr_arg = congrArg = let_val_congr
  Type: forall (b : a -> b), a = a' -> b a = b a'

  Nat.add_sub_cancel = Nat.add_sub_cancel_right
                     = Nat.add_sub_self_right
  Type: forall (a b : Nat), a + b - b = a

Size 2:
  mul_one = MulOneClass.mul_one
  Type: forall [MulOneClass M] (a : M), a * 1 = a

  mul_comm = CommMagma.mul_comm
  Type: forall [CommMagma G] (a b : G), a * b = b * a
\end{lstlisting}

The types shown are simplified pretty-printed representations of the
canonical forms; binder names and implicit arguments are artifacts of
the printer, not of the hash.
These duplicates arise from distinct sources.  The \texttt{by\_contra}
cluster contains four names for the same classical reasoning principle,
accumulated as Lean's standard library evolved.  The \texttt{mul\_one} /
\texttt{MulOneClass.mul\_one} pair reflects the algebraic hierarchy: the
root name and the typeclass-qualified name coexist as aliases.  The
\texttt{Nat.add\_sub\_cancel} triple shows the same arithmetic fact
named by three conventions.  In all cases, the content address correctly
identifies these as stating the same mathematical fact, modulo Level~1
canonicalization.

\paragraph{Coverage.}
All 57{,}613 theorems were hashed without error, demonstrating that
the pipeline (including the micro-batching workaround for
\texttt{whnf}; Appendix~\ref{app:engineering}) handles the full
transitive closure of a Mathlib module.  The current evaluation covers
a single module's environment at Level~1 only.  A systematic comparison
of Level~0 vs.\ Level~1 across multiple Mathlib entry points, and
evaluation in content-based (Merkle DAG) mode, remain as future work.

\simone{Quantitative evaluation across Mathlib: coverage, identical hashes at L0 vs L1, genuine duplicates, false positive rate, processing time. Produce a table across representative modules. CRITICAL.}



% =====================================================================
